%%%%%
%
\section{�chantillonnage et percentile par rapport a val\_{seuil} d'une variable sur une nouvelle base
de temps donn�e par un time\_{bnds} en tenant compte de valeurs invalides}

\itemstyle{Op\'erateur~:} \operfmt{PercentileInvEch\_tps\_bnds}
\index{PercentileInvEch\_tps\_bnds}

%%%%%OBJET%%%%
\itemstyle{Objet~:}
Cette fonction calcule, pour chaque instant $t$ d'une
base de temps $time_{out}$, 
une variable $X_{out}$, contenant le percentile associ� a un vecteur 
de valeurs de  $X_{in}$ entre les bornes donn�es par le $time_{bnds}$ associ� � $time_{out}$.
Le calcul n'est r�alis� que pour des valeurs < $val_{inv}$.
Aucune supposition n'est faite sur la r�gularit� de
l'�chantillonnage des 2 variables. Il est cependant n�cessaire
que $time_{in}$ (le temps de $X_{in}$) soit "� peu pr�s" tri� par ordre
croissant.


%%%%%ARGUMENTS%%%%
\bigskip
%% Dans ce fichier, les �ventuels accents sont cod�s
%% en ISO-8859-15 (aka latin9)
%% La ligne suivante donne trois exemples 
%% e accent aigu : �, e accent grave : �, e-dans-l'o : �
%%%%%ENTREES%%%%
\begin{tabarg}
\multicolumn{3}{l}{\titretab {Entr\'ees}}\\
\hline
$X_{in}$ & Vecteur, \texttt{float} & Variable d'origine\\
$toler_{mnq}$ & Coeff., \texttt{Int\_32} & Seuil de calcul du percentile (nb d'�chantillons). A chaque instant de la base de temps d'arriv�e, 
le percentile est calcul� si le nombre d'�chantillons dans l'intervalle de temps est strictement sup�rieur � $toler_{mnq}$. Sinon, 
le r�sultat est "valeur manquante". Si $toler_{mnq}$ = 0 le calcul est fait d�s qu'il y a au moins une valeur valide dans la p�riode\\
$time_{out}$ & Vecteur, \texttt{double} & Base de temps d'arriv�e (secondes)\\
$seuil$ & Coeff., \texttt{float} & Seuil du percentile (en pourcentage)\\
$val_{inv}$ & Coeff., \texttt{float} & valeur invalide pour laquelle le percentile n'est pas calcul� \\
$val_{remplac}$ & Coeff., \texttt{float} & valeur de remplacement quand la valeur atteint val\_{inv}  \\
\end{tabarg}

\medskip
%%%%%SORTIES%%%%
\begin{tabarg}
\multicolumn{3}{l}{\titretab {Sorties}}\\
\hline
$X_{out}$ & Vecteur, \texttt{float} & percentile de $X_{in}$ �chantillonn� sur
la base de $time_{out}$ (m�me unit� que $X_{in})$\\
\end{tabarg}



%%%%%FORMULE%%%%
\itemstyle{Formule}

%% PercentileEch_tps_bnds
Recherche du percentile pour la valeur $val_{seuil}$ inf�rieure a $val_{remplac}$ dans chaque �chantillon

Pour  i=$0$, Nbre  $Time_{bnds}$ :

	Pour chaque j entre ${tb[i][0]}$ et ${tb[i][1]}$ 

\[
X_{out}[i] =
\begin{cases}
\begin{tabular}{ll}
\texttt{ $X_{in}[j]$ telque $val_{seuil} \%$
 des valeurs sont sup�rieures}
&
\texttt{ si chaque $X_{in}[j]$ < $val_{inv}$}
\\
\texttt{$val_{replac}$} & \texttt{ sinon}
\end{tabular} 
\\
\end{cases}
\]


%%%%COMMENTAIRE%%%%
\itemstyle{Commentaire~:}
\textit{tb} est un tableau d'intervalles de temps d�termin�s � partir du vecteur temps $time_{out}$.

%%%%%AUTEUR%%%%
\itemstyle{Auteur~:}
CNRM/GMEI/TRAMM

%%%%%REFERENCE%%%%
\itemstyle{R\'ef\'erences~:}
N�ant.
